\documentclass[10pt,a4paper]{article}
\usepackage[T1]{fontenc}
\usepackage{geometry}
\geometry{margin=1.8cm,top=1.7cm,bottom=1.5cm,headheight=48pt,headsep=10pt}
\usepackage[scaled=0.9]{helvet}
\renewcommand{\familydefault}{\sfdefault}

\usepackage{tabularx} % in the preamble
\usepackage{longtable}
\usepackage{array}
\usepackage{hhline}
\usepackage{amssymb}
\usepackage{amsmath}
\usepackage{listings}
\usepackage{multirow}
\usepackage{fancyhdr}
\usepackage{wasysym}
\usepackage{adjustbox}
\usepackage{diagbox}
\usepackage{pdflscape}
\usepackage{float}
\usepackage{titlesec}
\usepackage{enumitem}
\usepackage{needspace}
\usepackage{multicol}
% A capture with no resolution of its own (a browser screenshot) prints at 150 pixels per inch,
% at most the text width: small crops are not blown up, and 13 px text prints at about 6 pt.
\ifdefined\pdfimageresolution\pdfimageresolution=150\fi
\BLOCK{ if record_layout }
% Evidence-record console output: breaks anywhere, even inside one long token.
\usepackage{fvextra}
\BLOCK{ endif }

\newcolumntype{R}[2]{%
    >{\adjustbox{angle=#1,lap=\width-(#2)}\bgroup}%
    l%
    <{\egroup}%
}
\newcommand*\rot{\multicolumn{1}{R{45}{-1em}}}

\makeatletter
\makeatother
\providecommand{\tabularnewline}{\\}

\usepackage[usenames,dvipsnames,svgnames,table,xcdraw]{xcolor} % use color
\usepackage{booktabs} % commands for table rules
\usepackage{tikz}
\usepackage{lastpage}

\usepackage[utf8]{inputenc}
\usepackage{hyperref}

\usepackage{pgfplots}
\DeclareUnicodeCharacter{2212}{−}
\usepgfplotslibrary{groupplots,dateplot}
\usetikzlibrary{patterns,shapes.arrows}
\pgfplotsset{compat=newest}

\usetikzlibrary{calc,shadings,patterns,tikzmark}

\newcommand\HatchedCell[4][0pt]{%
  \begin{tikzpicture}[overlay,remember picture]%
    \fill[#4] ( $ (pic cs:#2) + (0.0ex,1.9ex) $ ) rectangle ( $ (pic cs:#3) + (0.4ex,-#1*\baselineskip-.8ex) $ );
  \end{tikzpicture}%
}%

\newlength\colW \setlength\colW{0.096\columnwidth}

% Command hatch takes 5 args. Arg 1 = index into the HachtedCell command after table, arg 2 = width of cell in cols,
% arg 3 = the cell contents, args 4 & 5 are the cell side wall specifies, e.g. {|}{}
\newcommand*{\hatch}[5]{\multicolumn{#2}{#4!{\hspace*{-0.4pt}\tikzmark{start#1}} c !{\tikzmark{end#1}}#5}{#3}}
\newcommand*{\cellnormal}[4]{\multicolumn{#1}{#3 l #4}{#2}}

\newcommand*{\grayc}{\cellcolor{gray!25}}
\newcommand*{\darkgrayc}{\cellcolor{gray!65}}
\newcommand*{\errorc}{\cellcolor{red!50}}
\newcommand{\rom}{}

\newcommand*{\emptycheck}{\hspace{0.2cm}\makebox[0pt][l]{$\square$}{}\hspace{0.8cm}}
\newcommand*{\positivecheck}{\hspace{0.2cm}\makebox[0pt][l]{$\square$}\raisebox{.15ex}{\hspace{0.1em}$\checkmark$}\hspace{0.8cm}}
\newcommand*{\negativecheck}{\hspace{0.2cm}\makebox[0pt][l]{$\square$}{$\times$}\hspace{0.8cm}}
\definecolor{darkblue}{RGB}{0, 51, 102}
\definecolor{lightgray}{RGB}{240, 240, 240}
\newcommand{\traceflowtopleftlogo}{\VAR{top_left_logo}}
\newcommand{\traceflowtoprightlogo}{\VAR{top_right_logo}}
\newcommand{\traceflowstandardlogowidth}{4.0cm}
\newcommand{\traceflowwidelogowidth}{6.0cm}
% Right-header slot: expands to the logo image when one is set, otherwise to
% nothing. Used in every \fancyhead[R] (cover, portrait body and A3 landscape)
% so an omitted right logo never becomes an empty \includegraphics{}.
\BLOCK{ if top_right_logo }
\newcommand{\traceflowrightheaderstd}{\includegraphics[width=\traceflowstandardlogowidth,keepaspectratio]{\traceflowtoprightlogo}}
\newcommand{\traceflowrightheaderwide}{\includegraphics[width=\traceflowwidelogowidth,keepaspectratio]{\traceflowtoprightlogo}}
\BLOCK{ else }
\newcommand{\traceflowrightheaderstd}{}
\newcommand{\traceflowrightheaderwide}{}
\BLOCK{ endif }
\newcommand{\traceflowdocumentcode}{\VAR{document_code}}
\newcommand{\traceflowdocumenttype}{\VAR{document_type}}
\hypersetup{
  colorlinks=true,
  linkcolor=darkblue,
  urlcolor=blue,
  citecolor=darkblue
}
\urlstyle{same}

% Compact document spacing
\setlength{\parindent}{0pt}
\setlength{\parskip}{2pt}
\setlength{\LTpre}{0pt}
\setlength{\LTpost}{0pt}
\setlength{\textfloatsep}{8pt plus 2pt minus 2pt}
\setlength{\floatsep}{6pt plus 2pt minus 2pt}
\setlength{\intextsep}{6pt plus 2pt minus 2pt}
\renewcommand{\arraystretch}{1.15}
\setlist[itemize]{noitemsep, topsep=2pt, parsep=1pt, partopsep=0pt, leftmargin=1.5em}
\setlist[enumerate]{noitemsep, topsep=2pt, parsep=1pt, partopsep=0pt, leftmargin=1.5em}
\titleformat{\section}
  {\normalfont\large\bfseries\color{darkblue}}
  {}{0em}{}[\titlerule]
\titlespacing*{\section}{0pt}{1.2ex}{0.8ex}
\titleformat{\subsection}
  {\normalfont\normalsize\bfseries\color{darkblue}}
  {}{0em}{}
\titlespacing*{\subsection}{0pt}{0.8ex}{0.5ex}
\titleformat{\subsubsection}
  {\normalfont\normalsize\bfseries}
  {}{0em}{}
\titlespacing*{\subsubsection}{0pt}{0.5ex}{0.3ex}
% Console output: proportional spacing (fullflexible), small type, and breaks anywhere in
% a long line, each continuation marked with an arrow.
\lstset{
  basicstyle=\ttfamily\footnotesize,
  columns=fullflexible,
  keepspaces=true,
  breaklines=true,
  breakatwhitespace=false,
  postbreak=\mbox{\textcolor{gray}{$\hookrightarrow$}\space},
  showstringspaces=false,
  aboveskip=5pt,
  belowskip=5pt
}
% Let TeX stretch a line a little rather than push a long word past the margin.
\setlength{\emergencystretch}{3em}
\BLOCK{ if toc_depth is not none }
% Contents to the chosen depth. Headings are unnumbered, so the contents and the body
% agree and no number is glued to an ID; PDF bookmarks still go three levels deep.
\setcounter{tocdepth}{\VAR{toc_depth}}
\setcounter{secnumdepth}{0}
\hypersetup{bookmarksdepth=3}
% Room for three- and four-digit page numbers in the contents.
\makeatletter
\renewcommand{\@pnumwidth}{2.2em}
\makeatother
\BLOCK{ endif }

\begin{document}
\fancyhf{}
\fancyhead[L]{\includegraphics[width=\traceflowstandardlogowidth,keepaspectratio]{\traceflowtopleftlogo}}
\fancyhead[R]{\traceflowrightheaderstd}
\fancyhfoffset[L]{0pt}
\fancyhfoffset[R]{0pt}
\fancyfoot[L]{\footnotesize \VAR{footer_left}}
\fancyfoot[R]{\footnotesize Page \thepage\ of \pageref{LastPage}}
\renewcommand{\headrulewidth}{0.3pt}
\renewcommand{\footrulewidth}{0pt}
\pagestyle{fancy}
\setlength{\headwidth}{\textwidth}

\begin{center}
    \vspace*{0.4cm}
    {\LARGE \textbf{\VAR{report_title}}} \\[0.25em]
\BLOCK{ if cover_subtitle }
    {\large \VAR{cover_subtitle}} \\[0.25em]
\BLOCK{ endif }
    {\small \textbf{Generated:} \VAR{generated_date}}
\end{center}
\vspace{0.3cm}

\begin{tabularx}{\textwidth}{|X|X|}
\hline
\rowcolor{lightgray}
\textbf{Document ID} & \textbf{Document Type} \\
\hline
\traceflowdocumentcode{} & \traceflowdocumenttype{} \\
\hline
\end{tabularx}
\vspace{0.3cm}

\BLOCK{ if include_cover_page }
\BLOCK{ if cover_rows }
% A cover with no sign-off: the document is approved as part of another one.
\begin{table}[H]
\renewcommand{\arraystretch}{1.35}
\arrayrulecolor{gray} % Set the color of the horizontal and vertical lines to gray
\begin{tabular}{|>{\columncolor{gray!30}}m{0.3\linewidth}|m{0.6\linewidth}|}
\hline
\BLOCK{ for row in cover_rows }
\textbf{\VAR{row[0]}} & \VAR{row[1]} \\
\hline
\BLOCK{ endfor }
\end{tabular}
\end{table}
\BLOCK{ elif cover_signature_boxes }
\begin{table}[H]
\renewcommand{\arraystretch}{1.35}
\arrayrulecolor{gray} % Set the color of the horizontal and vertical lines to gray
\begin{tabular}{|>{\columncolor{gray!30}}m{0.45\linewidth}|m{0.45\linewidth}|}
\hline
\BLOCK{ if cover_version is not none }
\textbf{Version under test} & \VAR{cover_version} \\
\hline
\BLOCK{ endif }
\BLOCK{ if cover_engineer is not none }
\textbf{Engineer} & \VAR{cover_engineer} \\
\hline
\BLOCK{ endif }
\BLOCK{ if not cover_no_tester }
\textbf{Tester} & \\
\hline
\BLOCK{ endif }
\textbf{Test Date} & \VAR{cover_test_date} \\
\hline
\textbf{Initial Report Generation Date} & \VAR{generated_date} \\
\hline
\textbf{Result} & \VAR{cover_result} \\
\hline
\textbf{Approved Usage} \vspace*{1\baselineskip} & \VAR{cover_approved_usage} \\
\hline
\end{tabular}
\end{table}

\vspace{0.4cm}

\begin{minipage}{0.45\textwidth}
\centering
\fbox{\begin{minipage}[c][4.8cm][c]{0.92\linewidth}\mbox{}\end{minipage}}\\[0.6em]
\textbf{Lead Software Engineer}
\end{minipage}
\hfill
\begin{minipage}{0.45\textwidth}
\centering
\fbox{\begin{minipage}[c][4.8cm][c]{0.92\linewidth}\mbox{}\end{minipage}}\\[0.6em]
\textbf{Head of Core Lab}
\end{minipage}
\BLOCK{ else }
\begin{table}[H]
\renewcommand{\arraystretch}{1.35}
\arrayrulecolor{gray} % Set the color of the horizontal and vertical lines to gray
\begin{tabular}{|>{\columncolor{gray!30}}m{0.45\linewidth}|m{0.45\linewidth}|}
\hline
\BLOCK{ if cover_version is not none }
\textbf{Version under test} & \VAR{cover_version} \\
\hline
\BLOCK{ endif }
\BLOCK{ if cover_engineer is not none }
\textbf{Engineer} & \VAR{cover_engineer} \\
\hline
\BLOCK{ endif }
\BLOCK{ if not cover_no_tester }
\textbf{Tester} & \\
\hline
\BLOCK{ endif }
\textbf{Test Date} & \VAR{cover_test_date} \\
\hline
\textbf{Initial Report Generation Date} & \VAR{generated_date} \\
\hline
\textbf{Result} & \VAR{cover_result} \\
\hline
\textbf{Approved Usage} \vspace*{1\baselineskip} & \VAR{cover_approved_usage} \\
\hline
\end{tabular}
\end{table}

\vspace{0.4cm}


\BLOCK{ if not cover_single_signoff }
\begin{minipage}{0.45\textwidth}
\normalsize
\setlength{\baselineskip}{14pt}
\noindent \textbf{Manual tests carried out by:}

\noindent \rule{6cm}{0.4pt} \\
\noindent Signature

\noindent \rule{6cm}{0.4pt} \\
\noindent Name

\noindent \rule{6cm}{0.4pt} \\
\noindent Job Title

\noindent \rule{6cm}{0.4pt} \\
\noindent Date
\end{minipage}
\hfill
\BLOCK{ endif }
\begin{minipage}{0.45\textwidth}
\normalsize
\setlength{\baselineskip}{14pt}
\noindent \textbf{Validation pack result approved by:}
\BLOCK{ if cover_signature_meaning }

\noindent{\small Meaning of this signature: \VAR{cover_signature_meaning}\par}
\BLOCK{ endif }

\noindent \rule{6cm}{0.4pt} \\
\noindent Signature

\noindent \rule{6cm}{0.4pt} \\
\noindent Name

\noindent \rule{6cm}{0.4pt} \\
\noindent Job Title

\noindent \rule{6cm}{0.4pt} \\
\noindent Date
\end{minipage}
\BLOCK{ endif }
\BLOCK{ if cover_note }

\vspace{0.4cm}
\noindent{\small \VAR{cover_note}\par}
\BLOCK{ endif }
\BLOCK{ endif }

\BLOCK{ if include_cover_page or include_toc }
\newpage
\BLOCK{ endif }
\BLOCK{ if after_cover }
\VAR{after_cover}
\BLOCK{ endif }
\BLOCK{ if include_toc }
\tableofcontents
\BLOCK{ if toc_newpage }
\newpage
\BLOCK{ else }
\bigskip
\BLOCK{ endif }
\BLOCK{ endif }
